
\subsection*{6. Appendice~: \'Etude des $(b)$-faisceaux sur $P=\mathbb P^N_k$.}

\textbf{6.1.} Soient $k$ un corps alg\'ebriquement clos, et $P=\mathbb P^N_k$. Soient $F$ un faisceau coh\'erent sur $P$, \`a polyn\^ome de Hilbert
\[
   \chi(F(n))=\sum_{i=0}^{r} a_i \binom{n+i}{i},
\]
et
\[
   (c)=(c_0,\ldots,c_r)
\]
une suite d'entiers $\ge (a)$. Soient
\[
   a_{m,i}=\sum_{j=0}^{r-i}a_{j+i}\binom{m-1+j}{j}\quad\text{et}\quad c_{m,i}=\sum_{j=0}^{r-i}c_{j+i}\binom{m-1+j}{j}.
\]
Enfin, pour tout entier $q\ge1$, d\'esignons le plus grand faisceau coh\'erent de $F$ \`a support de dimension $\le q-1$, par $N_q$, et soit $F_q=F/N_q$.

\textbf{Lemme 6.2.} \emph{Sous les conditions de 6.1, consid\'erons une suite exacte}
\[
   0\longrightarrow I\longrightarrow {O}_P^{\oplus M}\longrightarrow F(m)\longrightarrow0
\]
\emph{avec $m\ge0$. Alors $I$ est $p$-r\'egulier avec \(p=P_r(c_{m,0},\ldots,c_{m,r})\) o\`u $P_r$ est le $r$-i\`eme $(b)$-polyn\^ome (1.9).}

En effet, ${O}_P$ est \'evidemment un $(0,\ldots,0,1)$-faisceau ($N+1$ termes), donc $I$ est un $(0,\ldots,0,M)$-faisceau d'apr\`es 1.6. Gr\^ace \`a 2.10, \(\chi(I(n))=M\binom{n+N}{N}-\sum_{i=0}^{r} a_{m,i}\binom{n+i}{i}\). Donc l'assertion r\'esulte de 1.11 et 1.10.

\textbf{Proposition 6.3.} \emph{Sous les conditions de 6.1, on suppose $F$ un $(b)$-faisceau. Alors on peut calculer par le moyen de polyn\^omes universels en les $a_i$, $b_i$ et $N$, des entiers $m, M, m_1$ et $M_1$ tels qu'on ait une suite exacte}
\[
   {O}_P(-m_1)^{\oplus M_1}\longrightarrow
   {O}_P(-m)^{\oplus M}\longrightarrow F\longrightarrow0.
\]
\emph{De plus, $N$ ne figure que dans le calcul de $M_1$.}

